\documentclass{article}
\usepackage[english]{babel}
\usepackage{amsmath}
\usepackage{amsthm}
\usepackage{amssymb}

\newtheorem{theorem}{Theorem}[section]
\newtheorem{lemma}[theorem]{Lemma}
\newenvironment{solution}{
\begin{proof}[Solution]}{
\end{proof}}

\begin{document}

% ============================================================
\section*{Section 2.3}

\subsection*{2.3.7}
Use the differentiation rules to find the derivative.

\subsubsection*{2.3.7 (c)}
\[
  f(z) = \frac{z^2 - 9}{iz^3 + 2z + \pi}
\]
\begin{solution}
  \begin{align*}
    f'(z)
    &= \frac{2z\left(iz^3 + 2z + \pi\right) - \left(3iz^2 + 2\right)\left(z^2 - 9\right)}{\left(iz^3 + 2z + \pi\right)^2} \quad \text{by rule (8)}
    \\ &= \frac{2iz^4 + 4z^2 + 2\pi z - \left[3iz^4 - 27iz^2 + 2z^2 - 18\right]}{\left(iz^3 + 2z + \pi\right)^2}
    \\ &= \frac{-iz^4 + z^2(2 + 27i) + 2\pi z + 18}{\left(iz^3 + 2z + \pi\right)^2}
  \end{align*}
\end{solution}

\subsection*{2.3.9}
Determine the points at which the function is not analytic.

\subsubsection*{2.3.9 (d)}
\[
  z^2\left(2z^2 - 3z + 1\right)^{-2}
\]
\begin{solution}
  Begin by finding the zeros of the expression inside the square in the denominator
  $(2z^2 - 3z + 1)$, which are the points where the function is non-analytic.

  The quadratic formula gives roots $\frac{1}{2}, 1$. Thus we can factor:
  \[
    \frac{z^2}{\left[\left(z - \frac{1}{2}\right)(z - 1)\right]^2}
    = \frac{z^2}{\left(z - \frac{1}{2}\right)^2(z - 1)^2}
  \]
  The denominator $= 0$ at the roots, so the function is non-analytic at
  $z = \frac{1}{2}$ and $z = 1$.
\end{solution}

\subsection*{2.3.10}
Let $f(z) = |z|^2$. Use the definition of the derivative to show that $f$ is
differentiable at $z = 0$ but not at any other point.
\begin{solution}
  Write the derivative from the definition and simplify algebraically exactly as given in the hint:
  \begin{align*}
    f'(z)
    &= \lim_{\Delta z \to 0} \frac{\lvert z_0 + \Delta z \rvert^2 - \lvert z_0 \rvert^2}{\Delta z}
    \\ &= \lim_{\Delta z \to 0} \left(\overline{z_0} + \overline{\Delta z} + z_0 \frac{\overline{\Delta z}}{\Delta z}\right)
    \\ &= \lim_{\Delta z \to 0} \left(\overline{z_0} + \overline{\Delta z}\right)
       + \lim_{\Delta z \to 0} \left(z_0 \frac{x - iy}{x + iy}\right)
  \end{align*}
  Note that $z_0$ is constant and $\Delta z \to 0$, so the first limit exists and
  goes to the constant $\overline{z_0}$.

  Take an approach direction for the second limit: $x = \theta y$. Note that
  $x \to 0 \implies \Delta z \to 0$.
  \[
    \lim_{x \to 0} z_0 \left(\frac{x - i\theta x}{x + i\theta x}\right)
    = z_0 \left(\frac{1 - \theta^2}{1 + \theta^2}\right)
  \]
  If $z_0 \neq 0$ the limit depends on the chosen approach angle $\theta$, so the
  derivative is different for different directions of $\Delta z$ and does not exist.

  If $z_0 = 0$ the limit does not depend on the approach angle and simply equals $0$.
  So the limit exists.
\end{solution}

\subsection*{2.3.16}
Let $f(z) = z^3 + 1$, $z_1 = \frac{-1 + \sqrt{3}i}{2}$, $z_2 = \frac{-1 - \sqrt{3}i}{2}$.
Show there is no $w$ on the segment from $z_1$ to $z_2$ with
$f(z_2) - f(z_1) = f'(w)(z_2 - z_1)$.
\begin{solution}
  We first divide both sides of the given equation by $z_2 - z_1$ and compute from
  the given $z_1, z_2$ values:
  \begin{align*}
    f'(w) &= \frac{f(z_2) - f(z_1)}{z_2 - z_1}
    \\ &= \frac{\dfrac{(-1 - \sqrt{3}i)^3}{8} - \dfrac{(-1 + \sqrt{3}i)^3}{8}}{z_2 - z_1}
    \\ &= \frac{(-1 - \sqrt{3}i)^3 - (-1 + \sqrt{3}i)^3}{-8\sqrt{3}i}
    \\ &= \frac{-6\sqrt{3}i + 6\sqrt{3}i}{-8\sqrt{3}i} = 0
    \quad \text{expanding out the cube and simplifying}
  \end{align*}
  So we need to find a $w$ satisfying $f'(w) = 0$ where $w$ is on the line between
  $z_1$ and $z_2$. We can parameterize this line as:
  \[
    w = -\frac{1}{2} + xi, \quad x \in \mathbb{R},\ -\frac{\sqrt{3}}{2} < x < \frac{\sqrt{3}}{2}
  \]
  Computing the derivative of $f$ and substituting:
  \[
    f(w) = w^3 + 1, \qquad f'(w) = 3w^2
  \]
  \[
    3\left(-\frac{1}{2} + xi\right)^2 = 3\left(\frac{1}{4} - xi - x^2\right)
  \]
  Set equal to zero and find values of $x$ satisfying the equation:
  \[
    \frac{3}{4} - 3x^2 - 3xi = 0
  \]
  Match real and imaginary parts:
  \begin{align*}
    \text{Im:} \quad & -3x = 0 \implies x = 0
    \\ \text{Re:} \quad & \frac{3}{4} - 3x^2 = 0 \implies \frac{3}{4} - 0 = \frac{3}{4} \neq 0
  \end{align*}
  As no values of $x$ satisfy this, no $w$ on the segment from $z_1$ to $z_2$
  satisfies $f(z_2) - f(z_1) = f'(w)(z_2 - z_1)$.
\end{solution}

% ============================================================
\section*{Section 2.4}

\subsection*{2.4.2}
Show that $h(z) = x^3 + 3xy^2 - 3x + i(y^3 + 3x^2y - 3y)$ is differentiable on
the coordinate axes but is nowhere analytic.
\begin{solution}
  Write the real and imaginary parts of $h$ as $u, v$ and compute partials:
  \begin{align*}
    u &= x^3 + 3xy^2 - 3x, & v &= y^3 + 3x^2y - 3y
    \\ u_x &= 3x^2 + 3y^2 - 3, & v_y &= 3y^2 + 3x^2 - 3
    \\ u_y &= 6xy, & v_x &= 6xy
  \end{align*}
  Check Cauchy-Riemann along the coordinate axes.

  Along the $y$ axis: $x = 0$, $y$ free:
  \begin{align*}
    u_x(0,y) &= 0 + 3y^2 - 3 = 3y^2 - 3 = v_y(0,y)
    \\ u_y(0,y) &= 0 = -v_x(0,y)
  \end{align*}
  Along the $x$ axis: $y = 0$, $x$ free:
  \begin{align*}
    u_x(x,0) &= 3x^2 + 0 - 3 = 3x^2 - 3 = v_y(x,0)
    \\ u_y(x,0) &= 0 = -v_x(x,0)
  \end{align*}
  Thus $h$ is differentiable along the coordinate axes.

  However, $h$ is non-analytic, as this fails on any open set. Take any deviation
  $\delta$ from the axes.

  Deviation from the $y$ axis:
  \[
    u_y(0 + \delta, y) = 6\delta y = -6\delta y = -v_x(0 + \delta, y)
  \]
  Only equal at $\delta = 0$ (for $y \neq 0$).

  Deviation from the $x$ axis:
  \[
    u_y(x, 0 + \delta) = 6x\delta = -6x\delta = -v_x(x, 0 + \delta)
  \]
  Only equal at $\delta = 0$ (for $x \neq 0$).

  Thus no nearby point off the axes is differentiable, and an open disk cannot be
  formed. Thus not analytic. As this holds for arbitrary $x, y$, the function is
  nowhere analytic.
\end{solution}

\subsection*{2.4.5}
Show that $f(z) = e^{x^2 - y^2}\left[\cos(2xy) + i\sin(2xy)\right]$ is entire, and
find its derivative.
\begin{solution}
  By Euler's formula, $\cos(2xy) + i\sin(2xy) = e^{2xyi}$, so:
  \begin{align*}
    f(z) &= e^{x^2 - y^2}\, e^{2xyi}
    \\ &= e^{x^2 + 2xyi - y^2}
    \\ &= e^{(x + yi)^2}
    \\ &= e^{z^2}
  \end{align*}
  $e^{z^2}$ is entire because it is a composition of entire functions: $z^2$ is a
  polynomial and $e^z$ is entire, and a composition of entire functions is entire by
  the chain rule.

  By the chain rule:
  \[
    \frac{d}{dz}\left[e^{z^2}\right] = e^{z^2} \cdot 2z
  \]
\end{solution}

\subsection*{2.4.6}
Show that in polar coordinates $(r, \theta)$ the Cauchy-Riemann equations become
\[
  \frac{\partial u}{\partial r} = \frac{1}{r}\frac{\partial v}{\partial \theta},
  \qquad
  \frac{\partial v}{\partial r} = -\frac{1}{r}\frac{\partial u}{\partial \theta}.
\]
\begin{solution}
  Write $z = re^{i\theta}$, $z_0 = r_0e^{i\theta_0}$, and compute
  $f'(z_0) = \lim_{z \to z_0} \frac{f(z) - f(z_0)}{z - z_0}$ along two paths as in the hint.

  Along the ray $\theta = \theta_0$: $z - z_0 = (r - r_0)e^{i\theta_0}$, so
  \begin{align*}
    f'(z_0)
    &= \lim_{r \to r_0} \frac{\left[u(r,\theta_0) - u(r_0,\theta_0)\right] + i\left[v(r,\theta_0) - v(r_0,\theta_0)\right]}{(r - r_0)e^{i\theta_0}}
    \\ &= e^{-i\theta_0}\left(u_r + iv_r\right)
  \end{align*}
  Along the circle $r = r_0$: $z - z_0 = r_0\left(e^{i\theta} - e^{i\theta_0}\right)$, and
  $\frac{e^{i\theta} - e^{i\theta_0}}{\theta - \theta_0} \to ie^{i\theta_0}$, so
  \begin{align*}
    f'(z_0)
    &= \lim_{\theta \to \theta_0} \frac{\left[u(r_0,\theta) - u(r_0,\theta_0)\right] + i\left[v(r_0,\theta) - v(r_0,\theta_0)\right]}{\theta - \theta_0}
       \cdot \frac{\theta - \theta_0}{r_0\left(e^{i\theta} - e^{i\theta_0}\right)}
    \\ &= \frac{u_\theta + iv_\theta}{ir_0e^{i\theta_0}}
    \\ &= e^{-i\theta_0}\,\frac{1}{r_0}\left(v_\theta - iu_\theta\right)
  \end{align*}
  Both are equal to $f'(z_0)$, so match real and imaginary parts:
  \begin{align*}
    \text{Re:} \quad & u_r = \frac{1}{r}v_\theta
    \\ \text{Im:} \quad & v_r = -\frac{1}{r}u_\theta
  \end{align*}
\end{solution}

\subsection*{2.4.10}
Show that if $f(z)$ is analytic and real-valued in a domain $D$, then $f(z)$ is
constant in $D$.
\begin{solution}
  Write $f(z)$ as its real and imaginary components:
  \[
    f(z) = u(z) + iv(z)
  \]
  If $f(z) \in \mathbb{R}$, then $v(z) = 0$.

  As $f$ is analytic, it satisfies Cauchy-Riemann, so
  \begin{align*}
    0 &= v_y = u_x
    \\ 0 &= v_x = -u_y
  \end{align*}
  So the partials of $u$ in $x$ and $y$ are both $0$.

  Integrate both components to find $u$:
  \begin{align*}
    u &= \int u_x \, dx = \int 0 \, dx = C_1 + g(y)
    \\ u &= \int u_y \, dy = \int 0 \, dy = C_2 + h(x)
  \end{align*}
  Setting equal:
  \[
    u(x + iy) = C_1 + g(y) = C_2 + h(x)
  \]
  Since this holds for all $x, y$, both sides must equal the same constant, so $g$
  and $h$ are constant. Thus $u = C$ and $f = C + 0i$ is constant.
\end{solution}

% ============================================================
\section*{Additional Problem}

\subsection*{Rotation}
Let $f(z)$ be analytic in a domain $D$. Prove that $e^{i\theta_0} f(z)$ is
analytic in $D$, where $\theta_0$ is a real constant.
\begin{proof}
  Let $f(z) = g(z) + ih(z)$ be analytic on $D$, with $g, h$ real.

  Expand $e^{i\theta} f(z)$ using Euler's identity:
  \begin{align*}
    e^{i\theta} f(z)
    &= (\cos\theta + i\sin\theta)\left[g(z) + ih(z)\right]
    \\ &= \underbrace{g\cos\theta - h\sin\theta}_{u(z)}
       + i\,\underbrace{\left(g\sin\theta + h\cos\theta\right)}_{v(z)}
  \end{align*}
  Check Cauchy-Riemann:
  \begin{align*}
    u_x &= g_x\cos\theta - h_x\sin\theta
    \\ v_y &= g_y\sin\theta + h_y\cos\theta
  \end{align*}
  Because $h_y = g_x$ and $g_y = -h_x$ (since $f$ is analytic):
  \[
    v_y = -h_x\sin\theta + g_x\cos\theta = u_x
  \]
  Same for the other Cauchy-Riemann equation:
  \begin{align*}
    u_y &= g_y\cos\theta - h_y\sin\theta
    \\ -v_x &= -g_x\sin\theta - h_x\cos\theta = -h_y\sin\theta + g_y\cos\theta = u_y
  \end{align*}
  The partials of $u, v$ are linear combinations of the partials of $g, h$, which
  are continuous, so they are continuous too.

  Thus $e^{i\theta} f(z)$ is analytic on $D$.
\end{proof}

% ============================================================
\section*{Section 2.4 (cont.)}

\subsection*{2.4.14}
Show that if the analytic function $w = f(z)$ maps a domain $D$ onto a portion
of a line, then $f$ must be constant throughout $D$.
\begin{proof}
  Write the line as $\left\{a + te^{i\alpha} : t \in \mathbb{R}\right\}$ for some
  fixed $a \in \mathbb{C}$ and $\alpha \in \mathbb{R}$. Then for each $z \in D$,
  \[
    f(z) = a + t(z)\,e^{i\alpha}, \qquad t(z) \in \mathbb{R}
  \]
  Define
  \[
    g(z) = e^{-i\alpha}\left(f(z) - a\right)
  \]
  We know that $g$ is analytic on $D$ because $a$ and $e^{-i\alpha}$ are constants,
  and adding or multiplying by constants does not change whether a function is
  analytic (previous problem). Then
  \[
    g(z) = e^{-i\alpha}\left(t(z)\,e^{i\alpha}\right) = t(z) \in \mathbb{R}
  \]
  An analytic function that is real-valued on a domain is constant (2.4.10), so
  $g(z) = C$ for some $C \in \mathbb{R}$. Thus
  \[
    f(z) = a + e^{i\alpha} g(z) = a + Ce^{i\alpha}
  \]
  which is constant throughout $D$.
\end{proof}

\end{document}
